\begin{figure}[tb]
\centering
\begin{tikzpicture}[scale=1.0]
 \draw[->] (-3,0)--(3.2,0) node[right]{$y_d$};
 \draw[->] (0,-2.1)--(0,2.25) node[above]{$y_e$};
 \draw[eigenvector] (0,0)--(2.15,1.25) node[above]{$z_\psi$: active flux};
 \draw[neutral axis] (0,0)--(-1.25,2.15) node[above left]{$z_0$: torque neutral};
 \draw[loss curve] (0,0) circle (1.55);
 \draw[->,gray] (1.0,0) arc[start angle=0,end angle=30,radius=1.0];
 \node[annotation] at (1.55,.25){loss-orthogonal rotation};
\end{tikzpicture}
\caption{After whitening, the active-flux coefficient selects one axis and
orthogonality uniquely supplies a torque-neutral complement. The quadrature
axis, perpendicular to this plane in 3D, remains unchanged.}
\label{fig:torque-axes}
\end{figure}
